\scanpage{94}
de $\pi_0$, $\pi'_0=\operatorname{Im}(\pi_1(\widetilde{\unsure{\mathcal{T}|\mathcal{C}}},x'_0)\unsure{-|\longrightarrow}\pi_1(\unsure{\mathcal{T}|\mathcal{C}},x_0))$, de sorte que
\[
\unsure{(47)}\qquad E_0\simeq\pi_0/\pi'_0;
\]
\deleted{$\unsure{\widetilde{\mathcal{T}}}$}
la condition d'existence de $\widetilde{\Theta}_0$ revient à dire
\[
\Theta(\pi'_0)\text{ conjugué à }\pi'_0\text{ dans }\pi_0
\]
\margin{en effet, il faudra que $\Theta(\pi'_0)$ soit $\operatorname{Im}(\pi_1(\widetilde{\unsure{\mathcal{T}|\mathcal{C}}},x'_1)\to\pi_1(\unsure{\mathcal{T}|\mathcal{C}},x_0))$, où $x'_1=\widetilde{\Theta}(x'_0)$, ou si $x'_1=g(x'_0)$, avec $g\in\pi_0$, l'Im en question est $\operatorname{int}(g)(\pi'_0)$.}
et les choix possibles pour remonter $\Theta$ en $\widetilde{\Theta}_0$, donc pour remonter $\Theta$ en $\widetilde{\Theta}$, correspondent aux éléments \inserted{$\unsure{x'_1}$} de $E_0$ \deleted{tels que} \inserted{de la} forme $g(x'_0)$ avec $g\in\pi_0$ satisfaisant
\[
(48)\qquad \operatorname{int}g(\pi'_0)=\Theta(\pi'_0)
\]
i.e. tels que
\[
g\in\operatorname{Transp}_{\pi_0}(\pi'_0,\Theta(\pi'_0))
\]
\margin{Cette condition ne dépend que de l'endom. \uline{ext.} défini par $\Theta$ sur $\pi'_0$ --- mais la façon de remonter $\Theta$ en $\widetilde{\Theta}$, $\unsure{=}$ \unsure{comme endom. ext., en dépend}, \unsure{aux de l'} \uline{\unsure{pas}} \unsure{un} endom. ord. \unsure{défini} par $\Theta$.}
\margin{\deleted{\unsure{Si} $\unsure{\widetilde{\mathcal{T}}}$ \unsure{correspond} \illegible{3} \unsure{on a} $\unsure{\Theta'_0}$ \illegible{3} $\unsure{\Theta}$ \illegible{2} $\unsure{\widetilde{\Theta}}$ \unsure{tel que} $\unsure{\exists\gamma\in\pi_0}$ \unsure{avec} $(48\,\mathrm{bis}):\ \unsure{\Theta}(\unsure{\pi'_0|x'_0})=\unsure{\gamma}(\unsure{\pi'_0|x'_0})$ \unsure{dans} $\pi'_0$.}}
donc les $\Theta$ possibles correspondent aux éléments de
\[
(49)\qquad \operatorname{Transp}_{\pi_0}(\pi'_0,\Theta(\pi'_0))/\pi'_0
\]
\margin{espace homogène à droite ($\pi'_0$ opère à droite sur le transporteur, via $g\longmapsto g\alpha$ pour $\alpha\in\pi_0$.}
\margin{Bien sûr, les façons de remonter $\Theta$: $\Theta'$ de $\mathcal{T}'\subset\widetilde{\mathcal{T}}$ correspondent aux $\widetilde{\Theta}$ qui rendent bien ainsi $\mathcal{T}'\subset\widetilde{\mathcal{T}}$ invariant.}

Ainsi, on voit que dans le cas \deleted{$\unsure{\mathcal{T}}$}
\[
\mathcal{T}=\widehat{\Pi}_1(X_{\unsure{\overline{k}}},I),\qquad I\subset X_{\unsure{k}}
\]
quand $\mathcal{T}'$ \deleted{\unsure{en premier}} est de la forme $\widehat{\Pi}_1(X'_{\unsure{\overline{k}}},I')$, où $X'$ est un rev. étale (fini) de $X$, \inserted{avec $I'\to I$}, que la condition qu'un $\Theta$ \inserted{de $\mathcal{T}$} se remonte en un \unsure{autom.} de $\mathcal{T}'$ est une condition algébrique des plus
\mspage{93}
